\section{Coefficientwise verification of the polynomial pullback}
\label{app:gauge-identities}

This appendix verifies that the compact design
\eqref{eq:B-compact}--\eqref{eq:C-compact} pulls back to the polynomial map
\eqref{eq:gauge-map}.  The calculation treats one generic coefficient \(g_k\)
at a time; it is not inferred from low-degree examples.

From \eqref{eq:Gpi} and \eqref{eq:beta},
\begin{equation}
 \beta(\Pi,S)
 =2\frac{g_2}{g_1}\Pi
  +\left(3\frac{g_3}{g_1}-1\right)\Pi S
  +\sum_{k=4}^N k\frac{g_k}{g_1}\Pi^kS^{k-2}.
 \label{eq:beta-expanded}
\end{equation}
On the source chart,
\[
 S=\frac{x}{t},\qquad \Pi=tq,\qquad Q=y+xq,
\]
and therefore, for every $k\geq4$,
\begin{equation}
 \Pi^kS^{k-2}=t^2x^{k-2}q^k.
 \label{eq:high-degree-B-monomial}
\end{equation}
Substitution now gives
\begin{align*}
 Q+\beta(\Pi,S)
 &=y+xq+2\frac{g_2}{g_1}tq
   +\left(3\frac{g_3}{g_1}-1\right)xq
   +\sum_{k=4}^N k\frac{g_k}{g_1}t^2x^{k-2}q^k\\
 &=y+3\frac{g_3}{g_1}xq+2\frac{g_2}{g_1}tq
   +\sum_{k=4}^N k\frac{g_k}{g_1}t^2x^{k-2}q^k,
\end{align*}
which is the displayed coordinate $B$.

For the third coordinate, write
\[
 C_*=\frac{2G_\Pi(S)}{g_1}-BS^2.
\]
Using $B=Q+\beta(\Pi,S)$ and the definition of $\beta$, we obtain
\begin{equation}
 C_*=2S-QS^2+\left(1-\frac{g_3}{g_1}\right)\Pi S^3
 -\sum_{k=4}^N(k-2)\frac{g_k}{g_1}\Pi^kS^k.
 \label{eq:C-separated}
\end{equation}
The terms with $k\geq4$ already have the desired source form, since
\begin{equation}
 \Pi^kS^k=(xq)^k.
 \label{eq:high-degree-C-monomial}
\end{equation}
It remains to simplify the cubic part.  Put $r=g_3/g_1$.  Then
\begin{align*}
 2S-QS^2+(1-r)\Pi S^3
 &=\frac{2xt-x^2y-rx^3q}{t^2}\\
 &=\frac{x(t+1)-rx^3q}{t^2},
\end{align*}
where the second line uses $xy=t-1$.  Since
\[
 q=t^2z+\frac1r y^2(1+3t),
\]
this becomes
\[
 -rx^3z+\frac{x\bigl(t+1-(t-1)^2(1+3t)\bigr)}{t^2}.
\]
The elementary identity
\begin{equation}
 t+1-(t-1)^2(1+3t)=t^2(5-3t)
 \label{eq:cubic-cancellation}
\end{equation}
therefore gives
\[
 2S-QS^2+(1-r)\Pi S^3=x(5-3t)-\frac{g_3}{g_1}x^3z.
\]
Together with \eqref{eq:C-separated} and
\eqref{eq:high-degree-C-monomial}, this is exactly the displayed coordinate
$C$.

We derived the identities after inverting $t$, but both sides are polynomials
in $K[x,y,z]$.  Localization at the nonzero polynomial $t$ is injective, so
the identities hold everywhere.  At fixed $\Pi$, their differential has
determinant
\[
 \det\frac{\partial(B,C)}{\partial(S,Q)}
 =-2\left(\frac{G_\Pi'(S)}{g_1}-BS\right)=-2D.
\]
The source-chart determinant is $D^{-1}$, and the chain rule gives
$\det DF_G=-2$.  This is the complete cancellation: each coefficient $g_k$
enters independently, so the same calculation works in every degree.
